\section{OSL-Core 0.2: mechanized finite control layer}
\label{sec:osl-v02}

\OSLCore 0.2 is the next finite implementation step.  It does not attempt to
mechanize the complete higher-order theory, but it moves several obligations
from prose into executable syntax and runtime objects.  Programs are parsed
into an explicit abstract syntax tree, checked before execution, and then
interpreted by a deterministic state machine.  Proof terms, temporal
contracts, transactions, finite-state sessions, bounded capabilities, neural
boundary proposals, and paired-trace metadata audits are represented directly.

\begin{statusbox}
The accompanying implementation contains 53 executable programs, 26 focused
automated tests, two paired-trace noninterference audits, an explicit parser and
static checker, and JSON schemas for packets, proofs, and capabilities.  All 53
programs match their declared acceptance or rejection status.  The safe audit
pair preserves the same public projection under two different secrets; the
leaky pair is rejected by the audit because public metadata changes with the
secret.  These are bounded executable checks, not a universal proof of the
full language.
\end{statusbox}

\subsection{Explicit abstract syntax and static checking}

Version 0.1 interpreted source lines directly.  Version 0.2 first constructs a
finite syntax tree
\[
  \operatorname{AST}(P)
  = (s_1,\ldots,s_n),
\]
where every statement stores its command kind, normalized arguments, source
line, and line number.  A declaration-oriented static pass verifies, before
runtime mutation, that:
\begin{itemize}[leftmargin=1.6em]
  \item lineages and agents are declared before use;
  \item operator source and target spaces exist;
  \item packets, proofs, actions, contracts, sessions, and secrets are named
  consistently;
  \item proof premises are already declared;
  \item capability effects and memory provenance classes are members of the
  finite language vocabulary;
  \item transactions are properly nested in the supported nonnested fragment;
  \item session participants and protocol references exist.
\end{itemize}

The static report is itself canonicalized and receipt-bound before execution.
This does not replace dynamic checking.  Matrix idempotence, proof freshness,
capability exhaustion, contract deadlines, and session-state transitions remain
runtime properties.

\begin{proposition}[Static pre-execution separation]
If a program fails the OSL-Core 0.2 static pass, then no program declaration or
runtime effect from that program is executed by the reference interpreter.
\end{proposition}

\begin{proof}
The interpreter calls the parser and static checker before entering the
statement-execution loop.  A static exception is returned before the runtime
step counter is advanced or any source statement is dispatched.  The only
state possibly present is the empty initial configuration.
\end{proof}

\subsection{First-class proof terms}

Proofs are now named runtime objects rather than only implicit conditions.  A
proof term has the finite form
\[
 p=
 \left(
 \begin{array}{l}
 \mathsf{name},\mathsf{kind},\mathsf{subject},\mathsf{agent},
 \mathsf{scope},\mathsf{memoryClass},\\
 \mathsf{sourceAgent},\mathsf{sourceLineage},
 \mathsf{premises},\mathsf{claims},\\
 d_{\mathrm{content}},d_{\mathrm{state}},t,
 \mathsf{trusted},d_p
 \end{array}
 \right).
\]
The stable content digest binds the proof to semantic content and original
provenance.  Composite use and ownership proofs additionally contain the
current mutable packet-state digest.  This makes them invalid after a relevant
state mutation.

Primitive proof constructors include:
\begin{lstlisting}[style=osl]
proof p_access access route_at_A by A
proof p_origin origin route_at_A source B lineage L_B
proof p_clear clear route_at_A for A
proof p_end endorsement route_at_A by A
proof p_auth authorization route_at_A by A for planning
proof p_obs observation event by A
proof p_cap capability externalAct by A
\end{lstlisting}

A composite planning proof is:
\begin{lstlisting}[style=osl]
proof p_use use route_at_A by A for planning
    from ["p_access","p_origin","p_clear","p_end","p_auth"]
\end{lstlisting}
The one-line concrete syntax is used in executable files; the line break above
is typographical.

An autobiographical ownership proof is:
\begin{lstlisting}[style=osl]
proof p_own ownership event by A as observed
    from ["p_access","p_origin","p_clear","p_observation"]
\end{lstlisting}

\begin{theorem}[Finite proof-state freshness]
Let $p$ be a composite use or ownership proof constructed for packet $c$ at
packet-state digest $d$.  If a later state mutation changes the canonical packet
digest to $d'\neq d$, then $p$ is rejected by every OSL-Core 0.2 instruction
requiring a fresh composite proof.
\end{theorem}

\begin{proof}
Composite proof validation requires equality between the stored
$d_{\mathrm{state}}(p)$ and the current packet digest.  The mutation changes the
latter to $d'$, while the immutable proof retains $d$.  The equality premise
therefore fails and execution reaches an explicit rejection.
\end{proof}

This is a finite revocation mechanism.  For example, a proof constructed before
an additional authorization or quarantine transition cannot silently authorize
a later plan.

\subsection{Proof-carrying planning example}

\lstinputlisting[style=osl,caption={A communicated observation is carried into planning by a named proof object.},label={lst:v02-proof-plan}]{examples/17_proof_planning.osl}

The proof path can be summarized as
\[
\begin{tikzcd}[column sep=small]
\Access_A(c)
  \arrow[r]
& \Origin_B(c)
  \arrow[r]
& \operatorname{Clear}_A(c)
  \arrow[r]
& K_A(c)
  \arrow[r]
& \Auth_A(c,\mathsf{planning})
  \arrow[r]
& \operatorname{UseProof}_A(c).
\end{tikzcd}
\]
The action separately consumes a bounded external-action capability.  Proof of
content use and authority to emit an effect remain distinct.

\subsection{Bounded and expiring capabilities}

A capability is represented as
\[
  \operatorname{Cap}
  (a,e,n,t_{\max},\rho,d),
\]
where $n\in\mathbb N\cup\{\infty\}$ is the remaining use budget and
$t_{\max}\in\mathbb N\cup\{\infty\}$ is the final valid runtime step.  The
surface syntax is:
\begin{lstlisting}[style=osl]
grant A effect externalAct uses 2 until 100
revoke A effect externalAct
\end{lstlisting}
Execution succeeds only if the capability exists, has not expired, and has
positive remaining use budget.  A successful execution decrements a finite
budget.

\begin{proposition}[Finite capability confinement]
In the OSL-Core 0.2 runtime, an action declaring effect $e$ cannot enter the
executed state unless the acting agent possesses a nonexpired capability for
$e$ with positive remaining budget.
\end{proposition}

\begin{proof}
The execution transition resolves the exact effect capability before mutating
the action.  Missing, expired, or exhausted capabilities raise a rejection.  On
success, the capability digest is bound into the action and a finite use count
is decremented.
\end{proof}

\lstinputlisting[style=osl,caption={A one-use tool capability is consumed by execution.},label={lst:v02-capability}]{examples/18_bounded_capability.osl}

\subsection{Temporal contracts}

The finite contract layer supports invariant templates and bounded response
obligations.  Examples are:
\begin{lstlisting}[style=osl]
contract no_foreign always no_foreign_observed
contract no_simulated always no_simulated_observed
contract no_quarantine_endorse always quarantine_blocks_endorsement
contract capability_required always capability_confinement
contract feedback_required after execute require feedback within 3
\end{lstlisting}

The feedback contract maintains a finite obligation map
\[
  \mathcal O_t:\mathsf{ActionName}\rightharpoonup\mathbb N,
\]
where $\mathcal O_t(u)$ is the due step for feedback on executed action $u$.
An execution inserts a deadline.  Matching feedback removes it.  A due step
passed without discharge, or an outstanding obligation at program termination,
causes rejection.

\begin{theorem}[Finite deadline-monitor soundness]
For every accepted OSL-Core 0.2 program containing a contract
\[
  \operatorname{execute}(u)
  \Rightarrow
  \Diamond_{\leq k}\operatorname{feedback}(u),
\]
every executed action occurring after contract declaration has a matching
feedback event no later than the monitor deadline.
\end{theorem}

\begin{proof}
The monitor inserts $u\mapsto t+k$ when the execute receipt is emitted.  A
matching feedback receipt removes $u$.  The runtime checks for overdue
obligations after every statement and treats every remaining obligation as a
violation at termination.  Hence any accepted execution has discharged every
inserted obligation within the declared bound.
\end{proof}

\lstinputlisting[style=osl,caption={A bounded response contract closes an action with feedback.},label={lst:v02-contract}]{examples/19_feedback_contract.osl}

\subsection{Transactional updates and rollback}

The transaction fragment is:
\begin{lstlisting}[style=osl]
begin_transaction repair
    # state-changing instructions
commit_transaction repair
\end{lstlisting}
or
\begin{lstlisting}[style=osl]
rollback_transaction repair
\end{lstlisting}

At transaction entry the finite runtime creates a deep checkpoint of the
canonical state.  Explicit rollback restores that checkpoint.  Any dynamic
rejection while the transaction is active also restores the checkpoint before
returning the rejection.  A rollback receipt records the checkpoint digest,
reason, begin step, and restored-state digest.

\begin{proposition}[Finite rollback atomicity]
Let $\Sigma_0$ be the runtime state immediately before transaction entry.  If a
transaction is explicitly or automatically rolled back, every program object
created after entry is absent from the restored application state, except for
the rollback audit receipt and the advanced audit step.
\end{proposition}

\begin{proof}
The runtime stores a deep copy of $\Sigma_0$ before transaction-body mutation.
Rollback replaces the active state with that copy and only then appends the
rollback receipt.  Objects created in the discarded branch are therefore not
members of the restored state maps.
\end{proof}

\lstinputlisting[style=osl,caption={A temporary autobiographical write is removed by explicit rollback.},label={lst:v02-rollback}]{examples/21_transaction_manual_rollback.osl}

Transactions are not nested in the finite language.  A production language
would require compositional transaction semantics, external-effect
compensation, durable storage, and concurrency control.

\subsection{Finite-state shared-workspace sessions}

A session specifies participants, an initial protocol state, and allowed
communication transitions:
\begin{lstlisting}[style=osl]
session CooperativeDiagnosis participants ["A","B"] initial waiting_for_A
session_transition CooperativeDiagnosis waiting_for_A
    communicate A B -> waiting_for_B
session_transition CooperativeDiagnosis waiting_for_B
    communicate B A -> ready
start_session CooperativeDiagnosis
\end{lstlisting}

A communication annotated with
\begin{lstlisting}[style=osl]
communicate map from A to B as map_at_B in CooperativeDiagnosis
\end{lstlisting}
can execute only when the protocol has a matching transition from its current
state.

\begin{theorem}[Finite session trace conformance]
Every accepted communication event labeled by session $S$ occurs along one of
$S$'s declared transitions, and the recorded session history is a path from the
initial state through those transitions.
\end{theorem}

\begin{proof}
The communication instruction queries the transition table at
\[
  (q,\mathsf{communicate},a,b),
\]
where $q$ is the current session state. Missing entries reject before packet
creation. Existing entries update the state to the declared target and append
the same edge to session history. Induction on the number of session-labeled
communications gives the result.
\end{proof}

\lstinputlisting[style=osl,caption={A two-agent evidence-exchange session enforces message order.},label={lst:v02-session}]{examples/24_session_protocol.osl}

This is a finite session-type monitor rather than a full asynchronous session
type system.  It does not yet prove deadlock freedom, liveness under scheduler
assumptions, or multiparty compatibility beyond the declared trace table.

\subsection{Neural boundary proposals}

The neural-to-symbolic interface does not convert a decoded activation directly
into trusted memory or action.  It creates a provisional packet:
\begin{lstlisting}[style=osl]
neural_propose latent_claim to A feature feature_441
    content "route-may-be-blocked" confidence 0.84
\end{lstlisting}
with status $\mathsf{provisional}$, provenance $\mathsf{unresolved}$, immediate
quarantine, and obligations
\[
  \{\mathsf{origin},\mathsf{lineage},\mathsf{task\_relevance}\}.
\]
Trusted boundary components may then issue attestations:
\begin{lstlisting}[style=osl]
attest origin latent_claim by A lineage L_A as p_origin
attest relevance latent_claim by A for planning as p_relevance
admit latent_claim by A using ["p_origin","p_relevance"]
\end{lstlisting}
Admission changes the packet to an inferred, self-bound, nonquarantined state.
It creates no direct-observation token.

\lstinputlisting[style=osl,caption={A latent proposal remains quarantined until trusted origin and relevance obligations are discharged.},label={lst:v02-neural}]{examples/22_neural_candidate_admission.osl}

\begin{boundarybox}
The attestation primitives are explicit trusted-boundary assumptions.  The
finite implementation does not infer causal origin from neural activations.  A
real neural boundary would need causal interventions, model-specific feature
identification, calibration, and adversarial validation.  OSL-Core 0.2 proves
only that an unvalidated proposal cannot silently bypass the declared
quarantine and admission path.
\end{boundarybox}

\subsection{Paired-trace noninterference audit}

Programs may register private secret values and publish public metadata.  The
public projection excludes secret values and retains only their digests.  The
paired auditor compares two accepted executions satisfying different registered
secret digests.  Let $\pi$ and $\pi'$ be the two traces and
$\operatorname{Pub}$ the declared public projection.  The tested property is
\[
  \operatorname{SecretDigest}(\pi)
  \neq
  \operatorname{SecretDigest}(\pi')
  \quad\Longrightarrow\quad
  \operatorname{Pub}(\pi)
  =
  \operatorname{Pub}(\pi').
\]

The safe pair uses trusted constant metadata:
\begin{lstlisting}[style=osl]
secret private_goal value "alpha"
metadata timestamp value 7 source trusted
\end{lstlisting}
and the paired run changes only the secret value.  The public projections are
identical.

The leaky pair publishes agent-controlled metadata derived from the secret:
\begin{lstlisting}[style=osl]
metadata covert from_secret private_goal source agent
\end{lstlisting}
The audit reports a public difference and returns failure.

\begin{proposition}[Paired audit correctness]
For two accepted finite runs, the OSL-Core 0.2 auditor reports
noninterference exactly when their canonical public projections are equal after
removing the derived public-digest field.
\end{proposition}

\begin{proof}
The auditor recursively compares every key, list element, type, and scalar value
of the two canonical projections.  It returns success precisely when the
difference set is empty.
\end{proof}

\begin{boundarybox}
This is a paired-execution test, not a HyperLTL proof over every possible
secret, context, schedule, message, or latent state.  It can falsify
noninterference for a pair and can verify equality for a declared pair.  It
cannot establish universal covert-channel absence.
\end{boundarybox}

\subsection{Mechanized finite guarantees and their limits}

The implemented core now checks the following properties directly:
\begin{enumerate}[leftmargin=1.7em]
  \item every accepted source file has a parsed AST and successful static report;
  \item direct observation is the only ordinary constructor of observation tokens;
  \item communicated, simulated, injected, and provisional neural contents do
  not become direct observations without the missing constructor;
  \item named composite proofs are packet-version sensitive;
  \item finite capability budgets and expiration steps are enforced at execution;
  \item declared finite temporal contracts are checked after every statement and
  at termination;
  \item transaction rejection restores the checkpointed application state;
  \item session-labeled messages follow the declared finite-state protocol;
  \item public paired traces can be compared for metadata noninterference;
  \item fixed source, backend, and environment produce deterministic snapshots
  and proof digests.
\end{enumerate}

These checks do not constitute a proof-assistant mechanization of preservation
and progress for the full language.  In particular, the parser and runtime are
ordinary Python programs, trusted attestations remain trusted assumptions,
external effects are modeled rather than sandboxed, and concurrency is absent.

\subsection{Expanded implementation map}

\begin{verbatim}
reference/
|-- osl_model.py              # canonical runtime and proof objects
|-- osl_parser.py             # AST parser and static checker
|-- osl_core.py               # deterministic operational semantics
|-- osl_check.py              # static-check CLI
|-- audit_noninterference.py  # paired-trace public audit
|-- run_examples.py           # 53 programs plus two audit pairs
|-- osl_core.ebnf             # concrete grammar
|-- packet.schema.json
|-- proof.schema.json
`-- capability.schema.json
examples/
|-- 01_... through 36_...     # accepted programs
`-- 91_... through 107_...   # expected rejections
tests/
|-- test_osl_core.py
`-- test_osl_core_v02.py
\end{verbatim}

\subsection{Next metatheoretic step}

The next formal milestone is to move the finite data model and operational rules
into a proof assistant.  The target suite is:
\begin{enumerate}[leftmargin=1.7em]
  \item preservation for the typed AST;
  \item progress to a value, next state, or declared rejection;
  \item linear provenance non-forging;
  \item effect and capability confinement;
  \item temporal-monitor soundness;
  \item transaction rollback refinement;
  \item session fidelity;
  \item compiler correctness for neural proposal lowering.
\end{enumerate}
Until then, the executable core should be interpreted as a transparent
reference semantics with adversarial tests, not as a small trusted kernel.
